\documentclass[10pt,a4paper]{amsart}
\usepackage[margin=19mm]{geometry}
\usepackage{amsmath,amssymb,mathtools}
\usepackage[scr=rsfs,cal=euler]{mathalfa}
\usepackage{xfrac}
\usepackage[unicode,hidelinks,bookmarks=false]{hyperref}
\newcommand{\ie}{i.e.\ }
\newcommand{\cf}{cf.\ }
\newcommand{\resp}{respectively\ }
\newcommand{\Subsection}{\subsection}
\providecommand{\nobreakdash}{\nobreak\hbox{-}}
\setlength{\emergencystretch}{2em}
\multlinegap0pt
\usepackage{amssymb}
\newtheorem{lemma}{Lemma}
\newtheorem{prop}{Proposition}
\newcommand{\supp}{\mathop{\mathrm{supp}}}
\title{Semiclassical shape resonances for magnetic Stark Hamiltonians}
\author{Kentaro Kameoka, Naoya Yoshida}
\date{Source: arXiv:2603.27196v1 (CC BY 4.0)}
\begin{document}
\maketitle

\noindent\emph{Source and licence.} Semiclassical shape resonances for magnetic Stark Hamiltonians. Version arXiv:2603.27196v1 (CC BY 4.0), distributed by arXiv under the Creative Commons Attribution 4.0 International licence, http://creativecommons.org/licenses/by/4.0/, as recorded in the arXiv licence metadata for this exact version and shown on the abstract page https://arxiv.org/abs/2603.27196v1. Full licence notice: \texttt{licenses/arxiv-2603.27196-CC-BY-4.0.txt}.\par\smallskip

\noindent\textit{Editorial scope.} Counted unit: the article's prop prop-projection (source0.tex lines 883--889). The article's complete original proof of it is delivered with it (source0.tex lines 892--923, one \texttt{proof} environment of 32 lines). No mathematics is written, repaired or re-derived: the statement and the proof are the article's own lines, in the article's own notation, and no retained line is substituted or re-flowed.  Retained uncounted as the source-local context the proof needs: lines 857--872.  Nothing was omitted from any retained span, so the package carries no omission record.  No retained line contains a citation, so the wrapper carries no bibliography and no bibliography entry.  No retained line contains a figure environment, an image-inclusion command or drawing code, so no image is delivered and the package declares no asset dependency.  Editorial formatting only, in addition: an AMS \texttt{amsart} preamble that restates the article's own declarations and macros used by the retained lines, namely the environments \texttt{lemma}, \texttt{prop} on separate counters, together with the standard packages those lines need.  The retained environments print in this document as Lemma 1, Proposition 1 in order of appearance, whereas the article prints the same items with the numbering of its own full document.  The complete native source of the article as distributed in the arXiv e-print for this exact version is bundled unchanged as \texttt{sources/197329/source0.tex}.
\begin{lemma}\label{lemma-cluster}
For small $h>0$, there exist $a_j(h)<b_j(h)<a_{j+1}(h)$ such that 
\[\left(\mathrm{Res}(P(h))\cup 
\sigma(P^{\mathrm{int}})\right)\cap ([a-\frac{\delta}{2}, b+\frac{\delta}{2}]-i[0, e^{-S_0/h}])
\subset \cup_{j=1}^{J(h)}\Omega_j(h),\]
where $\Omega_j(h)=[a_j(h), b_j(h)]-i[0, e^{-S_0/h}]$, 
$b_j-a_j\le Ch^{-2}e^{-S_0/h}$, $a_{j+1}-b_j\ge 2e^{-S_0/h}$, 
$a_1\in(a-\frac{2}{3}\delta, a-\frac{1}{3}\delta)$, $b_{J(h)}\in (b+\frac{1}{3}\delta,b+\frac{2}{3}\delta)$ and 
$\mathrm{Res}(P)\cap(([a_1-ch^2, a_1]-i[0, e^{-S_0/h}])
\cup ([b_{J(h)}, b_{J(h)}+ch^2]-i[0, e^{-S_0/h}]))
=\emptyset$.
Moreover, 
\[\|(1-\chi_0)(P_{\theta}-z)^{-1}\|\le C\alpha(h)^{-1},\mspace{7mu} z\in \partial \widetilde{\Omega}_j(h),\]
where $\widetilde{\Omega}_j(h)=[a_j(h)-e^{-S_0/h}, 
b_j(h)+e^{-S_0/h}]+i[-2e^{-S_0/h}, e^{-S_0/h}]$.
\end{lemma}

\begin{prop}\label{prop-projection}
For any $0<S<S_0$, 
\[ 
\Pi_j^{\theta}= \Pi_j^{\mathrm{int}}+\mathcal{O}(e^{-S/h})\,\,\,\,\,\mathrm{as} \,\,\,h\to 0
\]

\end{prop}

\begin{proof}
Since $z-P^{\mathrm{int}}$ is elliptic near $\supp (P_{\theta}-P^{\mathrm{int}})$, we have 
\begin{align*}
\|&(P_{\theta}-P^{\mathrm{int}})(z-P^{\mathrm{int}})^{-1}\chi_0\|\\
&\le C\|(z-P^{\mathrm{int}})(P_{\theta}-P^{\mathrm{int}})(z-P^{\mathrm{int}})^{-1}\chi_0\|_{L^2\to H_m^{-2}} \\
&= C \|[P^{\mathrm{int}}, P_{\theta}-P^{\mathrm{int}}](z-P^{\mathrm{int}})^{-1}\chi_0\|_{L^2\to H_m^{-2}} \\
&\le C+C e^{-S_0/h}\|(z-P^{\mathrm{int}})^{-1}\|
\end{align*}
where the last inequality follows from the Agmon estimate for $P^{\mathrm{int}}$ 
and the fact that $[P^{\mathrm{int}}, P_{\theta}-P^{\mathrm{int}}]$ has bounded coefficients. 
Then 
\begin{align*}
\|(P_{\theta}-P^{\mathrm{int}})(z-P^{\mathrm{int}})^{-1}\chi_0\|\le C
\end{align*}
when $\mathrm{dist}(z, \sigma(P^{\mathrm{int}}))\ge e^{-S_0/h}$. 

Since $\supp \chi_0\cap \supp (P_{\theta}-P^{\mathrm{int}})=\emptyset$, we have
\begin{align*}
 \Pi_j^{\theta}-\Pi_j^{\mathrm{int}}
=\frac{1}{2\pi i}\int_{\partial \widetilde{\Omega}_j}
(z-P_{\theta})^{-1}(1-\chi_0)(P_{\theta}-P^{\mathrm{int}})(z-P^{\mathrm{int}})^{-1}dz. 
\end{align*}
This and Lemma~\ref{lemma-cluster} imply 
\[
\| \left(\Pi_j^{\theta}- \Pi_j^{\mathrm{int}}\right) \chi_0\|\le C|\partial \widetilde{\Omega}_j| \alpha(h)^{-1} 
=\mathcal{O}(e^{-S/h}).
\] 
Finally, we have $\|\Pi_j^{\theta}(1-\chi_0)\|\le C|\partial \widetilde{\Omega}_j| \alpha(h)^{-1} 
=\mathcal{O}(e^{-S/h})$ by Lemma~\ref{lemma-cluster}, 
and $\|(1-\chi_0)\Pi_j^{\mathrm{int}}\|\le Ch^{-2}e^{-S_0/h}
=\mathcal{O}(e^{-S/h})$ by the Agmon estimate.
\end{proof}

\end{document}
